\documentclass[12pt,a4paper]{article}
\usepackage[utf8]{inputenc}
\usepackage[T1]{fontenc}
\usepackage{polski}
\usepackage{geometry}
\geometry{a4paper, margin=2.5cm}
\usepackage{graphicx}
\usepackage{hyperref}
\usepackage{float}
\usepackage{amsmath}
\usepackage{amssymb}
\usepackage{booktabs}
\usepackage{caption}
\usepackage{subcaption}

% Komenda pomocnicza do tworzenia "pustych miejsc" na zdjęcia
\newcommand{\placeholderimage}[2]{
    \begin{figure}[H]
        \centering
        \framebox(400,200){Miejsce na zdjęcie: #1}
        \caption{#2}
    \end{figure}
}

\title{\Huge Raport Techniczny Projektu:\\ \Large Wizja Komputerowa -- Detekcja i Klasyfikacja Znaków Drogowych}
\author{Jakub Jaszcz, Karol Malinowski, Jakub Śliwa, Filip Konfederak, Krystian Nowak}
\date{\today}

\begin{document}

\maketitle
\thispagestyle{empty}
\newpage

\tableofcontents
\newpage

\section{Wstęp (opis problemu, cel projektu)}

\subsection{Opis problemu}
Rozpoznawanie znaków drogowych (Traffic Sign Recognition, TSR) stanowi jeden z najważniejszych filarów rozwoju inteligentnych systemów transportowych (ITS) oraz zaawansowanych systemów wspomagania kierowcy (Advanced Driver Assistance Systems, ADAS). Znaki drogowe projektowane są w taki sposób, aby były łatwo rozpoznawalne dla człowieka -- charakteryzują się ścisłymi regułami dotyczącymi kształtu (np. koło, trójkąt, ośmiokąt) oraz koloru (np. dominacja czerwieni, niebieskiego, bieli i czerni). Niemniej jednak, dla systemów wizji komputerowej, ich poprawna detekcja i klasyfikacja w warunkach rzeczywistych jest zadaniem niezwykle skomplikowanym.

Systemy wizyjne operujące w rzeczywistym świecie muszą mierzyć się z szeregiem zniekształceń i anomalii środowiskowych. Należą do nich przede wszystkim:
\begin{itemize}
    \item \textbf{Zmienne warunki oświetleniowe:} Ostre światło słoneczne powodujące odblaski, głębokie cienie rzucane przez budynki, a także niedoświetlenie w warunkach nocnych.
    \item \textbf{Rozmycie w ruchu (Motion Blur):} Wynikające z wysokiej prędkości poruszającego się pojazdu w stosunku do czasu naświetlania klatki z kamery.
    \item \textbf{Przesłonięcia (Occlusions):} Znaki często są częściowo zasłonięte przez gałęzie drzew, infrastrukturę miejską, inne pojazdy czy nawet naklejki wandalistyczne.
    \item \textbf{Zniekształcenia perspektywiczne:} Kamera rzadko patrzy na znak idealnie na wprost. Wynikowa transformacja aficzna lub perspektywiczna zaburza idealne proporcje znaku.
    \item \textbf{Degradacja i warunki pogodowe:} Deszcz, śnieg, blaknięcie pigmentu od promieniowania UV.
\end{itemize}

Powyższe wyzwania sprawiają, że prosty szablonowy (template matching) algorytm rozpoznawania obrazów staje się bezużyteczny, wymuszając zastosowanie zaawansowanych algorytmów ekstrakcji cech oraz solidnych klasyfikatorów.

\subsection{Cel projektu}
Głównym celem niniejszego projektu było zaprojektowanie, zaimplementowanie oraz weryfikacja skuteczności systemu wizyjnego do detekcji, segmentacji i klasyfikacji wybranych klas znaków drogowych. W celu głębokiego zrozumienia mechanizmów wizji komputerowej, świadomie zrezygnowano z wykorzystania potężnych, pretrenowanych konwolucyjnych sieci neuronowych (np. YOLO czy ResNet) na rzecz klasycznych, interpretowalnych algorytmów przetwarzania obrazów. 

Projekt realizuje potok (pipeline) składający się z:
\begin{enumerate}
    \item Wstępnego przetwarzania obrazu i filtracji szumów.
    \item Detekcji obszarów zainteresowania (ROI) przy pomocy analizy geometrycznej oraz stabilnych regionów ekstremalnych (MSER).
    \item Wyodrębnienia niezmienniczych cech wektorowych za pomocą deskryptora HOG.
    \item Klasyfikacji znaku maszyną wektorów nośnych (SVM).
\end{enumerate}


\section{Opis zastosowanych metod}
Wdrożony system dzieli się na dwa główne moduły: lokalizację znaku na pełnym zdjęciu (segmentacja) oraz rozpoznanie odnalezionego znaku (ekstrakcja cech i klasyfikacja). W ramach segmentacji przetestowano i zaimplementowano dwa równoległe podejścia.

\subsection{Podejście A: Klasyczna analiza krawędzi i geometrii}
Metoda ta opiera się na założeniu, że znaki drogowe posiadają silnie zarysowane, standardowe obrysy geometryczne. Przetwarzanie obrazu w tym podejściu wymaga kilku operacji matematycznych na macierzy pikseli:

\subsubsection{Konwersja do skali szarości i rozmycie Gaussa}
Pierwszym krokiem jest redukcja informacji z trzech kanałów barw (RGB) do pojedynczego kanału jasności (Luminance). Następnie obraz poddawany jest filtrowaniu dolnoprzepustowemu za pomocą rozmycia Gaussa. Maska konwolucyjna (kernel) o wymiarach $3\times3$ przesuwa się po obrazie, uśredniając wartości pikseli zgodnie z rozkładem normalnym. Zabieg ten jest krytyczny, ponieważ eliminuje szumy wysokiej częstotliwości -- np. teksturę asfaltu czy liście, które mogłyby wygenerować miliony niepotrzebnych krawędzi w kolejnych etapach.

\placeholderimage{Zestawienie obrazu oryginalnego i po zastosowaniu Gaussian Blur}{Porównanie obrazu przed i po wygładzaniu Gaussem. Wstaw tutaj screeny działania Twojego algorytmu.}

\subsubsection{Detekcja Krawędzi Canny'ego i Operacje Morfologiczne}
Algorytm Canny'ego to wieloetapowy proces wykrywania krawędzi. Oblicza on gradienty jasności obrazu (zazwyczaj operatorem Sobela), a następnie wykonuje tłumienie wartości niemaksymalnych (Non-maximum suppression), zostawiając tylko piksele o najostrzejszym gradiencie. Na koniec stosuje się progowanie z histerezą (w projekcie użyto progów 50 i 150).

Z uwagi na to, że obiekty w świecie rzeczywistym nie zawsze mają ciągłe krawędzie (np. z powodu odblasków świetlnych), na wynikowy obraz binarny nałożono operację morfologicznej \textbf{dylatacji}. Wykorzystując strukturę elementarną (maskę jedynek $3\times3$), dylatacja powiększa obszary białe (krawędzie), efektywnie łącząc rozerwane obrysy znaków w spójne figury zamknięte.

\placeholderimage{Obraz po Canny Edge Detection i Dylatacji}{Widok binarnych krawędzi znaku odnalezionych przez algorytm. Widoczne zamknięte obrysy.}

\subsubsection{Detekcja konturów i aproksymacja Ramer-Douglas-Peucker}
Na uzyskanym obrazie krawędziowym algorytm wyszukuje zamknięte topologicznie krzywe (kontury). Aby odrzucić nieregularne, naturalne obiekty, wykorzystano funkcję aproksymacji wielokątów. Algorytm Ramer-Douglas-Peucker redukuje liczbę wierzchołków krzywej na podstawie zadanej tolerancji (w kodzie ustalono margines błędu na 8\% obwodu konturu). Akceptowane są obiekty posiadające co najmniej 3 wierzchołki (trójkąty, czworokąty, ośmiokąty, wielokąty przybliżające okrąg) z zachowaniem proporcji boków (Aspect Ratio) od 0.4 do 2.0.

\subsubsection{Weryfikacja barwą w przestrzeni HSV}
Kształt wielokąta nie gwarantuje obecności znaku -- może to być okno budynku lub koło samochodu. Dlatego dla zidentyfikowanych "kandydatów" przeprowadzana jest ostateczna weryfikacja. Obraz konwertowany jest do walcowej przestrzeni barw HSV (Hue, Saturation, Value). Przestrzeń ta separuje informacje o jasności (Value) od samej barwy (Hue), co znacznie uodparnia system na zmianę nasłonecznienia. Nałożono maski dla koloru czerwonego i niebieskiego. Kandydat zostaje zaakceptowany, jeśli piksele pożądanej barwy stanowią minimum 10\% jego powierzchni.

\placeholderimage{Maska HSV na wyciętym znaku}{Po lewej wycięty kandydat w RGB, po prawej białe piksele oznaczające spełnienie maski dla koloru czerwonego w przestrzeni HSV.}


\subsection{Podejście B: Detekcja Plam MSER (Maximally Stable Extremal Regions)}
Alternatywną metodą zaimplementowaną w projekcie jest ekstrakcja obszarów MSER. Algorytm ten iteracyjnie zmienia próg jasności obrazu (od 0 do 255) i obserwuje, jak zachowują się "plamy" (spójne obszary binarne). Jeżeli dany region nie zmienia znacząco swojego kształtu i pola powierzchni przez szeroki zakres progów, zostaje uznany za "maksymalnie stabilny ekstremalny region".

Znaki drogowe, ze względu na duży kontrast piktogramów w stosunku do własnego tła (np. czarna liczba 50 na białym tle z czerwoną obwódką) są idealnymi kandydatami dla detektora MSER.
Parametry algorytmu ograniczono do wykrywania plam o polu powierzchni między 100 a 50000 pikseli. Podobnie jak w pierwszym podejściu, wynikowe "plamy" przefiltrowano weryfikatorem proporcji (wymuszającym kształt zbliżony do kwadratu) oraz walidatorem barw HSV (wymagając 15\% nasycenia koloru znaku).

\placeholderimage{Detekcja MSER na zdjęciu wejściowym}{Plamy stabilne odnalezione przez detektor MSER obrysowane na zdjęciu testowym.}


\subsection{Ekstrakcja Cech -- Histogram Skierowanych Gradientów (HOG)}
Kluczowym elementem odpowiadającym za skuteczność systemu wizyjnego jest deskryptor HOG (Histogram of Oriented Gradients). Zamienia on macierz pikseli wejściowego, wykadrowanego znaku w spłaszczony wektor matematyczny, niosący informacje o kształcie.

Algorytm HOG operuje według następujących kroków:
\begin{enumerate}
    \item \textbf{Obliczanie gradientów:} Dla każdego piksela wyliczana jest amplituda i kąt gradientu zmian jasności. Największe wartości powstają na ostrych krawędziach piktogramu znaku.
    \item \textbf{Podział na komórki (Cells):} Obraz znormalizowany do wymiaru $64\times64$ pikseli dzielony jest na komórki o wielkości $8\times8$ pikseli.
    \item \textbf{Tworzenie histogramów:} Dla każdej komórki tworzony jest 9-przedziałowy (9 bins) histogram orientacji gradientów w zakresie 0-180 stopni. Każdy piksel głosuje na dany przedział z siłą proporcjonalną do swojej amplitudy.
    \item \textbf{Normalizacja w blokach (Blocks):} Komórki są grupowane w bloki o wielkości $16\times16$ px (zawierające po 4 komórki), nakładające się na siebie z krokiem 8 px. Normalizacja bloku (zazwyczaj metodą normy L2) uodparnia ostateczny wektor na globalne zacienienia znaku.
\end{enumerate}
W ten sposób obraz znaku zamieniany jest na wysokowymiarowy wektor liczbowy pozbawiony nieistotnych danych (jak dokładny kolor pojedynczych pikseli).


\subsection{Klasyfikacja maszyną wektorów nośnych (SVM)}
Wygenerowane wektory HOG przekazywane są do klasyfikatora Support Vector Machine (SVM). Celem SVM jest odnalezienie w wielowymiarowej przestrzeni takich hiperpłaszczyzn, które najlepiej separują wektory należące do różnych klas.

W projekcie wykorzystano jądro liniowe (\texttt{kernel='linear'}). Jest to uzasadnione matematycznie -- w przypadku deskryptorów HOG uzyskujemy przestrzeń o ogromnej liczbie wymiarów (dla parametrów z projektu jest to wektor rzędu kilku tysięcy wartości), w której dane zazwyczaj stają się liniowo separowalne. Użycie bardziej skomplikowanych jąder (jak Radial Basis Function -- RBF) mogłoby prowadzić do nadmiernego dopasowania (overfittingu) oraz drastycznie zwiększyło czas treningu i predykcji. Margines tolerancji ustawiono na standardowym poziomie $C=1.0$.


\section{Opis wykorzystanych danych}
Zastosowanie i ewaluację zaprezentowanych metod przeprowadzono w oparciu o zestaw danych GTSRB (German Traffic Sign Recognition Benchmark). Baza ta jest uważana za standard (benchmark) dla problematyki klasyfikacji znaków i została opracowana na poczet międzynarodowego konkursu w 2011 roku.

Oryginalny zbiór jest niezwykle bogaty i zróżnicowany:
\begin{itemize}
    \item Zawiera ponad 50 000 zdjęć przypisanych do 43 klas znaków.
    \item Zdjęcia cechują się dużą wariancją wielkości: od miniaturowych $15\times15$ pikseli do dużych, ostrych kadrów $250\times250$ pikseli.
    \item Baza udostępnia koordynaty precyzyjnych ramek ROI dla każdego znaku.
\end{itemize}

\subsection{Wybrane klasy i filtracja zbioru}
Ze względu na obszerność bazy oraz ograniczone zasoby obliczeniowe do treningu eksperymentalnych metod klasycznych, zawężono przestrzeń problemu do detekcji 5 charakterystycznych klas znaków, zróżnicowanych pod kątem kształtu i koloru:
\begin{enumerate}
    \item \textbf{Klasa 2:} Ograniczenie prędkości do 50 km/h (okrągły, obwódka czerwona, czarny piktogram).
    \item \textbf{Klasa 12:} Droga z pierwszeństwem (romb, żółty środek, biała obwódka).
    \item \textbf{Klasa 14:} Znak STOP (ośmiokąt, czerwone tło, biały napis).
    \item \textbf{Klasa 17:} Zakaz wjazdu (okrągły, czerwone tło, biały prostokąt).
    \item \textbf{Klasa 35:} Nakaz jazdy na wprost (okrągły, niebieskie tło, biała strzałka).
\end{enumerate}

\placeholderimage{Przykłady zdjęć dla 5 wybranych klas z bazy GTSRB}{Kolaż zdjęć prezentujący próbkę każdej z pięciu wykorzystywanych w projekcie klas. Zwróć uwagę na znaczne różnice w oświetleniu.}

\subsection{Proces wstępnej obróbki (Preprocessing)}
Dedykowany skrypt przygotowujący dane poddał bazę następującym manipulacjom:
\begin{itemize}
    \item \textbf{Filtrowanie i selekcja:} Wczytanie plików \texttt{Train.csv} oraz \texttt{Test.csv} i wyselekcjonowanie tylko wierszy pasujących do 5 zdefiniowanych klas. Skrypt oddzielił zbiór treningowy (służący do dopasowania hiperpłaszczyzn SVM) od zbioru testowego (na którym weryfikowana jest skuteczność modelu).
    \item \textbf{Kadrowanie (Cropping):} Znaki wycięto z oryginalnych zdjęć z minimalnym tłem na podstawie dokładnych koordynatów obszaru zainteresowania (ROI.X1, ROI.Y1, ROI.X2, ROI.Y2).
    \item \textbf{Skalowanie (Resizing):} Interpolacja każdego wyciętego znaku do znormalizowanej rozdzielczości $64\times64$ pikseli. Skalowanie podyktowane jest bezwzględnym wymaganiem algorytmu HOG.
\end{itemize}


\section{Przeprowadzone eksperymenty i analiza wyników}

\subsection{Ewaluacja klasyfikacji modelu HOG + SVM}
Na precyzyjnie wyciętych i wykadrowanych znakach ze zbioru testowego uruchomiono potok klasyfikatora HOG+SVM w celu określenia jego ogólnej "zdolności rozumienia" piktogramów na znakach drogowych. 

Dla uzyskanych predykcji sporządzono raport klasyfikacji.
\textit{Poniżej znajduje się przygotowane miejsce na wstawienie tabeli z logów konsoli Twojego skryptu ewaluacyjnego:}

\begin{table}[H]
    \centering
    \begin{tabular}{l c c c c}
        \toprule
        \textbf{Klasa Znaku} & \textbf{Precision} & \textbf{Recall} & \textbf{F1-Score} & \textbf{Liczba próbek (Support)} \\
        \midrule
        2 (50 km/h)          & [wstaw] & [wstaw] & [wstaw] & [wstaw] \\
        12 (Pierwszeństwo)   & [wstaw] & [wstaw] & [wstaw] & [wstaw] \\
        14 (STOP)            & [wstaw] & [wstaw] & [wstaw] & [wstaw] \\
        17 (Zakaz Wjazdu)    & [wstaw] & [wstaw] & [wstaw] & [wstaw] \\
        35 (Nakaz na wprost) & [wstaw] & [wstaw] & [wstaw] & [wstaw] \\
        \midrule
        \textbf{Accuracy (Dokładność Ogólna):} & \multicolumn{4}{c}{\textbf{[WSTAW WARTOŚĆ \%]}} \\
        \bottomrule
    \end{tabular}
    \caption{Szczegółowy raport skuteczności klasyfikacji modelu HOG+SVM dla 5 wybranych klas.}
\end{table}

\subsection{Analiza Macierzy Pomyłek (Confusion Matrix)}
Macierz pomyłek pozwala na doskonałą wizualizację tego, z jakimi klasami model myli się najczęściej. Idealny model reprezentowany jest przez mocno nasyconą wartościami przekątną główną.

\placeholderimage{Macierz Pomyłek (Confusion Matrix)}{Wygenerowana macierz pomyłek. Oś X to wartości przewidywane (Predicted), oś Y to wartości prawdziwe (True Label).}

\textit{[Omówienie macierzy pomyłek: na podstawie wstawionego wykresu przeanalizuj tutaj krótko, czy model częściej myli np. okrągły znak "50" ze znakiem zakazu wjazdu ze względu na ten sam zarys, czy może występują rzadkie pomylenia rombu ustąp pierwszeństwa ze znakiem STOP.]}

\subsection{Analiza Przypadków Błędnych (Error Analysis)}
Projekt wykorzystuje unikalny podsystem "Error Trackingu". Skrypt klasyfikacyjny każdorazowo tworzy i wypełnia folder \texttt{errors/}, kopiując do niego wszystkie zdjęcia na których model dokonał błędnej predykcji z opisowym formatem \texttt{True\_[x]\_Pred\_[y]\_[nazwa].png}. Przegląd tego folderu ukazuje słabe strony klasycznego podejścia HOG+SVM. 

Dominującymi problemami skutkującymi błędem modelu były:
\begin{enumerate}
    \item \textbf{Silne przesłonięcie i wandalizm:} Zdjęcia, na których znak obklejony jest wlepkami kibicowskimi zmienia całkowicie układ gradientów (HOG), przez co klasyfikator SVM widzi w nim losowy zlepek wektorów.
    \item \textbf{Skrajne ocienienie:} Zdjęcia w lesie z promieniami słońca przecinającymi znak w połowie generowały mocną sztuczną krawędź poziomą/pionową podczas uderzania gradientu na krawędź światłocienia.
\end{enumerate}

\placeholderimage{Zestawienie 3-4 błędnie sklasyfikowanych obrazków z folderu errors}{Przykłady zdjęć, które oszukały model HOG+SVM. Po lewej ekstremalne warunki świetlne, po prawej zniszczenia mechaniczne i wandalizm.}

\subsection{Omówienie wpływu modyfikacji parametrów detekcji i segmentacji}
W trakcie prac projektowych testowano różne wartości parametrów dla algorytmów geometrycznych, obserwując ich bezpośredni wpływ na odrzuty znaków w pierwszej fazie potoku:
\begin{itemize}
    \item \textbf{Rozmiar filtra Gaussa:} Zmniejszenie maski poniżej $3\times3$ powodowało wybuch ilości fałszywych krawędzi Canny'ego dla chropowatych tekstur tła.
    \item \textbf{Progi histerezy Canny'ego:} Obniżenie dolnego progu powodowało dołączenie słabo kontrastowych "szumów" do krawędzi obiektów, co czasami wręcz utrudniało domknięcie konturu znaku drogowego na tle mocno zabudowanego miasta.
    \item \textbf{Limit powierzchni MSER:} Podniesienie minimalnego limitu z 100 na wyższe wartości, skutkowało niezauważeniem znaków w oddali przez detektor MSER, z kolei jego zbytnie obniżenie generowało tysiące malutkich "plamek" z tekstury jezdni i szumu wideo, drastycznie spowalniając wykonanie programu.
\end{itemize}


\section{Wnioski i podsumowanie}

\subsection{Co działało poprawnie i zalety rozwiązania}
\begin{itemize}
    \item Kombinacja \textbf{deskryptora HOG oraz klasyfikatora SVM} udowodniła, że dla poprawnie odciętego z tła obiektu, statyczne podejście wciąż osiąga wybitne wyniki rozpoznawalności. HOG jest silnie odporny na równomierne ściemnienie obrazu lub drobną utratę jakości.
    \item Klasyfikacja za pomocą SVM przebiega ekstremalnie szybko -- trening modelu na macierzy cech to ułamek czasu, jakiego wymagałoby strojenie wag sieci konwolucyjnej, a predykcja jest niemalże natychmiastowa, co sprzyjałoby systemom osadzonym (Embedded Systems).
    \item Użycie masek w przestrzeni barw \textbf{HSV} było strzałem w dziesiątkę jako ostateczny bezpiecznik dla szalenie generujących obszary detektorów takich jak MSER. Filtrowanie barwą ucięło znaczącą ilość "False Positives".
\end{itemize}

\subsection{Co nie działało idealnie -- wyzwania i ograniczenia}
\begin{itemize}
    \item Podejście klasyczne silnie ufa rygorystycznym ramom proporcji (Aspect Ratio). Znaki sfotografowane z boku okna, a nie centralnie z przedniej szyby powodują zmianę konturu koła w ostry kształt zbliżony do elipsy. Wymusza to drastyczne poluzowanie rygoru proporcji i kątów dla geometrii krawędzi w OpenCV.
    \item Dylatacja ratuje przerwane po detektorze Canny'ego kontury, ale niebezpiecznie blisko rosnące wokół znaku gałęzie powodują, że proces ten "skleja" znak z drzewem. Wynikowy kontur staje się nieregularnym kształtem ameby i zostaje brutalnie odrzucony jako znak we wczesnej fazie odrzutu przed klasyfikacją.
\end{itemize}

\subsection{Rekomendacje i możliwe usprawnienia}
Biorąc pod uwagę gwałtowny postęp sztucznej inteligencji, docelowym usprawnieniem architektury oprogramowania w przyszłości byłoby wdrożenie dwuetapowej sieci głębokiej lub sieci typu Single-Shot Detector (np. \textbf{YOLOv8} bądź \textbf{SSD}). 

Rozwiązania oparte na Konwolucyjnych Sieciach Neuronowych (CNN) w mniejszym stopniu opierają się na ręcznie "zakodowanych" heurystykach (jak polecenie odrzucenia konkretnego koloru HSV), lecz same inferują cechy nieliniowe z setek tysięcy próbek. Zastosowanie pretrenowanej głębokiej sieci neuronowej znacznie poprawiłoby skuteczność detekcji (ramkowania) znaków na bardzo zaszumionym miejskim tle, częściowo ignorując naturalne przerwania piktogramu powstałe przez warunki atmosferyczne czy bliki świetlne.


\newpage
\section{Bibliografia}
\begin{enumerate}
    \item Zbiór danych GTSRB (German Traffic Sign Recognition Benchmark) w serwisie Kaggle: \url{https://www.kaggle.com/datasets/meowmeowmeowmeowmeow/gtsrb-german-traffic-sign}
    \item Oficjalna dokumentacja biblioteki OpenCV (Open Source Computer Vision Library): \url{https://docs.opencv.org/master/}
    \item Oficjalna dokumentacja biblioteki Scikit-Learn (Machine Learning in Python): \url{https://scikit-learn.org/stable/}
    \item Przewodnik użytkownika Scikit-Learn -- Maszyny Wektorów Nośnych (SVM): \url{https://scikit-learn.org/stable/modules/svm.html}
    \item Oficjalna dokumentacja biblioteki Pandas (użytej do obróbki plików CSV): \url{https://pandas.pydata.org/docs/}
\end{enumerate}

\end{document}
